\chapter{Aktuelle KI-Tools im Überblick}\label{ch:K03_Tools}

\begin{myremark}
Die KI-Landschaft entwickelt sich rasend schnell. In diesem Kapitel erhalten Sie
einen strukturierten Überblick über aktuelle Werkzeuge in den Bereichen Sprache,
Bild, Video und Audio – mit Stärken, Schwächen und konkreten Anwendungsbeispielen.
\end{myremark}

% ─────────────────────────────────────────────────────────────────────────────
\section{Grosse Sprachmodelle (LLMs) im Vergleich}

Die folgende Tabelle gibt einen Überblick über die bekanntesten LLMs (Stand Frühjahr
2025). Da sich die Modelle laufend weiterentwickeln, empfiehlt es sich, die aktuellen
Benchmarks auf \url{https://lmarena.ai} nachzuschlagen.

\begin{table}[H]
\centering
\caption{Vergleich aktueller LLMs (vereinfacht, Stand 2025)}
\label{tab:llm-vergleich}
\small
\begin{tabular}{p{2.2cm}p{2cm}p{3.5cm}p{3.5cm}p{2cm}}
    \toprule
    \textbf{Modell}     & \textbf{Anbieter} & \textbf{Stärken}                                         & \textbf{Schwächen}                              & \textbf{Zugang}      \\
    \midrule
    GPT-4o              & OpenAI            & Allrounder, multimodal (Text, Bild, Audio), Coding        & Kostenpflichtig (Vollversion), Datenschutz (USA)& Web, API             \\
    \midrule
    Claude 3.5 Sonnet   & Anthropic         & Langer Kontext, sicherheits-betont, starkes Schreiben     & Kein Bildgenerator, US-Unternehmen              & Web, API             \\
    \midrule
    Gemini 1.5 Pro      & Google            & Sehr langer Kontext (1\,M Tokens), Google-Integration     & Qualität teils inkonsistent                     & Web, API             \\
    \midrule
    Llama 3             & Meta              & Open Source, lokal betreibbar                            & Erfordert technisches Setup                     & Lokal / Hugging Face \\
    \midrule
    Mistral             & Mistral AI (FR)   & Europäischer Anbieter, Open Source möglich                & Kleiner als GPT-4, weniger Features             & Web, API             \\
    \midrule
    Brian               & CH-spezifisch     & Datenschutzkonform, für Schulen geeignet                 & Eingeschränkter Funktionsumfang                 & Schullizenz          \\
    \midrule
    Fobizz              & Fobizz GmbH (DE)  & DSGVO-konform, Bildungskontext, diverse Tools             & Lizenzkosten für Schulen                        & Schullizenz          \\
    \bottomrule
\end{tabular}
\end{table}

\begin{myexercise}[– Vergleich selbst testen]
Stellen Sie denselben Prompt an mindestens zwei verschiedene KI-Tools (z.\,B.
Fobizz / Brian und ein weiteres Ihrer Wahl). Verwenden Sie folgenden Prompt:

\begin{center}
    \textit{Erkläre mir das Konzept des maschinellen Lernens in 5 Sätzen für
    eine 16-jährige Person, die noch keine Vorkenntnisse hat.}
\end{center}

Vergleichen Sie die Antworten anhand der Kriterien: Verständlichkeit,
Korrektheit, Länge und Stil. Welches Modell hat am besten abgeschnitten?
Begründen Sie Ihre Antwort.
\end{myexercise}

\begin{myexercise}[– Modell-Rätsel]
\textbf{Welches Modell bin ich?} – Lesen Sie die folgenden Beschreibungen und
ordnen Sie jeder Beschreibung ein Modell aus \cref{tab:llm-vergleich} zu:

\begin{enumerate}
    \item \enquote{Ich bin vollständig quelloffen und kann auf Ihrem eigenen
          Computer betrieben werden.}
    \item \enquote{Ich komme aus Frankreich und bin stolz darauf, ein europäischer
          Anbieter zu sein.}
    \item \enquote{Ich kann einen Roman von 700 Seiten auf einmal lesen und
          beantworten.}
    \item \enquote{Mein Unternehmen legt besonderen Wert auf Sicherheit und
          ethischen Umgang mit KI.}
\end{enumerate}
\end{myexercise}

% ─────────────────────────────────────────────────────────────────────────────
\section{KI-Agenten}

\begin{mydefinition}
Ein \textbf{KI-Agent} ist ein System, das nicht nur Fragen beantwortet, sondern
selbstständig \textbf{Aufgaben planen und ausführen} kann. Dazu nutzt es
verschiedene \textbf{Tools} (Werkzeuge) wie z.\,B. Websuche, Code-Ausführung
oder das Versenden von E-Mails.

Bekannte Agentensysteme: OpenAI Operator, Claude Computer Use,
Microsoft Copilot Agent, AutoGPT, LangChain.
\end{mydefinition}

\begin{myexample}
Ein Agent erhält den Auftrag: \enquote{Buche mir den günstigsten Flug von Zürich
nach Berlin nächsten Freitag und sende mir eine Bestätigungs-E-Mail.}

Der Agent könnte:
\begin{enumerate}
    \item Eine Websuche nach Flügen durchführen.
    \item Die Ergebnisse vergleichen.
    \item Den günstigsten Flug auswählen und den Buchungsprozess starten.
    \item Eine Bestätigungs-E-Mail generieren und senden.
\end{enumerate}
\end{myexercise}

\begin{myremark}
Weitere Details zu Agenten finden Sie in \cref{ch:K05_Agenten}.
\end{myremark}

% ─────────────────────────────────────────────────────────────────────────────
\section{Bildgeneratoren}

\begin{table}[H]
\centering
\caption{Vergleich aktueller KI-Bildgeneratoren (Stand 2025)}
\label{tab:bildgen-vergleich}
\small
\begin{tabular}{p{2.5cm}p{2cm}p{3.5cm}p{3.5cm}p{2cm}}
    \toprule
    \textbf{Tool}      & \textbf{Anbieter}    & \textbf{Stärken}                            & \textbf{Schwächen}                     & \textbf{Zugang}   \\
    \midrule
    DALL-E 3           & OpenAI               & Prompt-Integration in ChatGPT, hohe Qualität & Kostenpflichtig                        & ChatGPT Plus      \\
    \midrule
    Midjourney         & Midjourney           & Sehr hohe künstlerische Qualität             & Nur via Discord, Abo nötig             & Discord           \\
    \midrule
    Stable Diffusion   & Stability AI         & Open Source, lokal betreibbar                & Setup-Aufwand, Qualität variiert       & Lokal / Web       \\
    \midrule
    Adobe Firefly      & Adobe                & Lizenzfrei (kommerzielle Nutzung), sicher    & Erfordert Adobe-Konto                  & Web               \\
    \midrule
    Imagine (Meta)     & Meta                 & Kostenlos, in WhatsApp / Instagram integriert & Eingeschränkte Kontrolle, Datenschutz & Meta-Apps         \\
    \bottomrule
\end{tabular}
\end{table}

\begin{myexercise}[– Bildgenerator ausprobieren]
Nutzen Sie Adobe Firefly (\url{https://firefly.adobe.com}) – dieser Dienst
verwendet nur lizenzfreie Bilder und eignet sich für den Schulkontext.

Generieren Sie mit folgendem Prompt ein Bild:
\begin{center}
    \textit{A Swiss mountain landscape at sunrise, digital art style,
    high quality, no text}
\end{center}

Variieren Sie anschliessend den Prompt: Fügen Sie einen Stil hinzu
(z.\,B. \enquote{oil painting}, \enquote{pixel art}) oder ändern Sie
die Szene. Was fällt Ihnen auf?
\end{myexercise}

\begin{mychallenge}[Midjourney – Künstlerische Bildgenerierung]
\begin{myattention}
\textbf{Optional – Eigenverantwortung:} Für Midjourney ist ein Discord-Account
und ein kostenpflichtiges Abo erforderlich. Nutzung auf eigene Verantwortung.
\end{myattention}

Erkunden Sie Midjourney (\url{https://midjourney.com}). Erstellen Sie ein Bild
im Stil eines bekannten Künstlers (z.\,B. Van Gogh, Banksy) und beschreiben Sie,
welche Prompt-Elemente den grössten Einfluss auf das Ergebnis hatten.

Diskutieren Sie: Ist die Verwendung von Künstlerstilen durch KI ethisch
vertretbar?
\end{mychallenge}

% ─────────────────────────────────────────────────────────────────────────────
\section{Videogeneratoren}

\begin{table}[H]
\centering
\caption{Vergleich aktueller KI-Videogeneratoren (Stand 2025)}
\label{tab:videogen-vergleich}
\small
\begin{tabular}{p{2.5cm}p{2cm}p{3.5cm}p{3.5cm}p{2cm}}
    \toprule
    \textbf{Tool}   & \textbf{Anbieter} & \textbf{Stärken}                               & \textbf{Schwächen}                           & \textbf{Zugang}    \\
    \midrule
    Sora            & OpenAI            & Sehr realistische kurze Videos, hohe Qualität   & Nur für Plus/Pro-Nutzer, USA                 & ChatGPT Plus/Pro   \\
    \midrule
    Luma Dream Machine & Luma AI        & Kostenloser Einstieg, photorealistische Videos  & Begrenzte kostenlose Generierungen           & Web (\url{lumalabs.ai}) \\
    \midrule
    Gemini (Veo 2)  & Google           & Lange Videoclips, hohe Qualität                 & Eingeschränkter Zugang                       & Google AI          \\
    \midrule
    Runway ML       & Runway           & Professionell, viele Editing-Features           & Teuer, Lernaufwand                           & Web               \\
    \bottomrule
\end{tabular}
\end{table}

\begin{mychallenge}[Videogenerator testen]
\begin{myattention}
\textbf{Optional – Eigenverantwortung:} Videogeneratoren sind kommerzielle Dienste.
Nutzung auf eigene Verantwortung; datenschutzkonforme Alternativen bevorzugen.
\end{myattention}

Besuchen Sie Luma Dream Machine (\url{https://lumalabs.ai/dream-machine}).
Generieren Sie ein kurzes Video (5 Sekunden) mit folgendem Prompt:
\begin{center}
    \textit{A cat walking on a snowy Swiss mountain path, cinematic style}
\end{center}

Reflektieren Sie:
\begin{itemize}
    \item Was beeindruckt Sie? Was wirkt noch \enquote{falsch}?
    \item Welche ethischen Probleme entstehen durch realistische Videogenerierung?
    \item Kennen Sie Beispiele, wo KI-generierte Videos missbraucht wurden?
\end{itemize}
\end{mychallenge}

% ─────────────────────────────────────────────────────────────────────────────
\section{Audio-Generatoren}

\begin{table}[H]
\centering
\caption{Aktuelle KI-Audio-Generatoren (Stand 2025)}
\label{tab:audiogen-vergleich}
\small
\begin{tabular}{p{2.5cm}p{2.5cm}p{3.5cm}p{3.5cm}}
    \toprule
    \textbf{Tool}    & \textbf{Anbieter}   & \textbf{Stärken}                              & \textbf{Schwächen}                         \\
    \midrule
    ElevenLabs       & ElevenLabs          & Stimmenkloning, sehr natürliche Sprache        & Missbrauchspotenzial (Deepfakes)           \\
    \midrule
    Suno             & Suno AI             & Vollständige Songs aus Textprompts             & Urheberrechtliche Fragen                   \\
    \midrule
    Udio             & Udio                & Hochwertige Musik, viele Stile                 & Eingeschränkte kostenlose Nutzung          \\
    \midrule
    Whisper          & OpenAI              & Transkription (Sprache → Text), Open Source    & Nur Transkription, kein Sprachgenerator    \\
    \bottomrule
\end{tabular}
\end{table}

\begin{myremark}
\textbf{Audio-Deepfakes} sind eine ernsthafte Gefahr: Mit wenigen Sekunden
Audiomaterial kann eine Stimme geklont werden. Dies wird für Betrug
(\enquote{Enkeltrick 2.0}) und Desinformation missbraucht.
\end{myremark}

\begin{myexercise}[– Deepfake-Erkennung]
Suchen Sie im Internet nach öffentlichen Beispielen von Audio- oder Video-Deepfakes
(z.\,B. auf der Website \url{https://www.media.mit.edu/projects/detect-fakes/overview/}).

\begin{itemize}
    \item Können Sie den Deepfake von der echten Aufnahme unterscheiden?
    \item Welche Merkmale verraten einen Deepfake?
    \item Welche gesellschaftlichen Konsequenzen hätte eine Welt, in der
          man Audio und Video nicht mehr vertrauen kann?
\end{itemize}
\end{myexercise}
